\documentclass[11pt]{article}
\usepackage[utf8]{inputenc}
\usepackage[T1]{fontenc}
\usepackage[spanish]{babel}
\usepackage{amsmath, amssymb, amsthm, mathtools}
\usepackage{braket}
\usepackage{dsfont}
\usepackage{graphicx}
\usepackage{float}
\usepackage[margin=2.5cm]{geometry}
\usepackage{hyperref}
\hypersetup{hidelinks}
\emergencystretch=2em
\allowdisplaybreaks

\newcommand{\ii}{\mathrm{i}}
\newcommand{\ee}{\mathrm{e}}
\newcommand{\Id}{\mathds{1}}

\title{Estado coherente fon\'onico en bloque polarit\'onico}
\author{Ricardo Del Rosal Gardu\~no \\ \small Doctorado en F\'isica, BUAP}
\date{Notas de doctorado}

\begin{document}
\maketitle
\tableofcontents
\bigskip

\section{Estado inicial y objetivo}

Se considera un bloque polarit\'onico fijo del modelo efectivo de
Holstein--Cummings. A diferencia del caso con un n\'umero fon\'onico definido,
el modo vibracional se prepara ahora en un estado coherente desplazado
\begin{equation}
  \boxed{
  \ket{\tilde\beta}
  =
  \sum_{m=0}^{\infty}\tilde c_m\ket m,
  \qquad
  \tilde c_m
  =
  \ee^{-M/2}\frac{\tilde\beta^m}{\sqrt{m!}},
  \qquad
  M=|\tilde\beta|^2.
  }
\end{equation}


La distribuci\'on fon\'onica es de Poisson:
\begin{equation}
  \boxed{
  |\tilde c_m|^2
  =
  \ee^{-M}\frac{M^m}{m!}.
  }
\end{equation}
El estado coherente contiene una superposici\'on de muchos valores de $m$.
Cada componente evoluciona con su propia frecuencia de Rabi efectiva, por lo
que la suma produce colapsos y resurgimientos de las observables.

El manuscrito se centra en el bloque $n=9$ y toma
\begin{equation}
  \boxed{
  \tilde\beta=3,
  \qquad
  M=9,
  \qquad
  \ket{\tilde\Psi(0)}=\ket{-,9}\ket{\tilde\beta}.
  }
\end{equation}
La preparaci\'on polarit\'onica inferior equivale a
\begin{equation}
  \boxed{
  d_9^+=0,
  \qquad
  d_9^-=1.
  }
\end{equation}

Se estudiar\'an dos niveles de descripci\'on. Primero se linealiza la frecuencia
de Rabi alrededor del n\'umero medio $M$. Esta aproximaci\'on permite expresar
la evoluci\'on mediante estados coherentes rotados y obtener anal\'iticamente el
tiempo de revival. Despu\'es se conserva la dependencia exacta con $m$ dentro
del hamiltoniano efectivo. Esa suma completa recupera la forma de los revivals,
su ensanchamiento y la reducci\'on de su altura.

\section{Par\'ametros del bloque \texorpdfstring{$n=9$}{n=9}}

Para un bloque $n$ y un n\'umero fon\'onico cercano a $M$, se definen
\begin{equation}
  \delta_n=2g\sqrt n-\nu,
  \qquad
  \Omega_{nM}^{\mathrm{ef}}
  =
  \sqrt{\delta_n^2+\nu^2\lambda^2M}.
\end{equation}
La derivada respecto al n\'umero fon\'onico, evaluada en $M$, es
\begin{equation}
  \boxed{
  \Omega_{nM}^{\mathrm{ef}\prime}
  =
  \left.
  \frac{\partial\Omega_{nm}^{\mathrm{ef}}}{\partial m}
  \right|_{m=M}
  =
  \frac{\nu^2\lambda^2}{2\Omega_{nM}^{\mathrm{ef}}}.
  }
\end{equation}
La frecuencia linealizada se escribe como
\begin{equation}
  \Omega_{nm}^{\mathrm{ef}}
  \simeq
  \varphi_{nM}+m\Omega_{nM}^{\mathrm{ef}\prime},
\end{equation}
donde
\begin{equation}
  \boxed{
  \varphi_{nM}
  =
  \Omega_{nM}^{\mathrm{ef}}
  -M\Omega_{nM}^{\mathrm{ef}\prime}.
  }
\end{equation}
Los coeficientes de mezcla se aproximan por su valor en el centro de la
distribuci\'on:
\begin{equation}
  C=C_{nM}
  =
  \sqrt{
  \frac{\Omega_{nM}^{\mathrm{ef}}+\delta_n}
  {2\Omega_{nM}^{\mathrm{ef}}}
  },
  \qquad
  S=S_{nM}
  =
  \sqrt{
  \frac{\Omega_{nM}^{\mathrm{ef}}-\delta_n}
  {2\Omega_{nM}^{\mathrm{ef}}}
  }.
\end{equation}
Por construcci\'on,
\begin{equation}
  C^2+S^2=1,
  \qquad
  C^2-S^2
  =
  \frac{\delta_n}{\Omega_{nM}^{\mathrm{ef}}},
  \qquad
  2CS
  =
  \frac{\nu\lambda\sqrt M}{\Omega_{nM}^{\mathrm{ef}}}.
\end{equation}

Para $n=M=9$ se obtiene
\begin{equation}
  \boxed{
  A_{99}
  =
  9\Omega+\frac{\nu\lambda^2}{4}+\frac{17}{2}\nu,
  \qquad
  \delta_9=6g-\nu.
  }
\end{equation}
La frecuencia efectiva y su derivada son
\begin{equation}
  \boxed{
  \Omega_{99}^{\mathrm{ef}}
  =
  \sqrt{\delta_9^2+9\nu^2\lambda^2}
  =
  \sqrt{(6g-\nu)^2+9\nu^2\lambda^2},
  }
\end{equation}
\begin{equation}
  \boxed{
  \Omega_{99}^{\mathrm{ef}\prime}
  =
  \frac{\nu^2\lambda^2}{2\Omega_{99}^{\mathrm{ef}}},
  \qquad
  \varphi_{99}
  =
  \Omega_{99}^{\mathrm{ef}}
  -9\Omega_{99}^{\mathrm{ef}\prime}.
  }
\end{equation}

\section{Estado aproximado como dos paquetes coherentes}

La aproximaci\'on lineal obtenida en la nota de din\'amica da, para un bloque
$n$ fijo,
\begin{align}
  \ket{\tilde\Psi(t)}
  &\simeq
  \ee^{-\ii A_{nM}t}
  \Bigg[
  (d_n^+C-d_n^-S)
  \ee^{-\ii\varphi_{nM}t/2}
  \ket{\tilde\beta\ee^{-\ii\Omega_{nM}^{\mathrm{ef}\prime}t/2}}
  \notag\\
  &\qquad\otimes
  \left(
  C\ee^{-\ii\Omega_{nM}^{\mathrm{ef}\prime}t/2}\ket{+,n}
  -S\ket{-,n}
  \right)
  \notag\\
  &\quad+
  (d_n^+S+d_n^-C)
  \ee^{+\ii\varphi_{nM}t/2}
  \ket{\tilde\beta\ee^{+\ii\Omega_{nM}^{\mathrm{ef}\prime}t/2}}
  \notag\\
  &\qquad\otimes
  \left(
  S\ee^{+\ii\Omega_{nM}^{\mathrm{ef}\prime}t/2}\ket{+,n}
  +C\ket{-,n}
  \right)
  \Bigg].
\end{align}
El estado vibracional inicial se divide en dos paquetes coherentes. Sus
amplitudes en el espacio de fase son
\begin{equation}
  \tilde\beta_-(t)
  =
  \tilde\beta\ee^{-\ii\Omega_{nM}^{\mathrm{ef}\prime}t/2},
  \qquad
  \tilde\beta_+(t)
  =
  \tilde\beta\ee^{+\ii\Omega_{nM}^{\mathrm{ef}\prime}t/2}.
\end{equation}
Los paquetes giran en sentidos opuestos. Su separaci\'on angular es
$\Omega_{nM}^{\mathrm{ef}\prime}t$. Cuando se separan, disminuye su traslape y
se pierde la interferencia entre las ramas polarit\'onicas. Cuando vuelven a
coincidir, reaparece la interferencia y ocurre un revival.

Para $n=M=9$, $d_9^+=0$ y $d_9^-=1$, el estado se reduce a
\begin{align}
  \ket{\tilde\Psi(t)}
  &\simeq
  \ee^{-\ii A_{99}t}
  \Bigg[
  -S\ee^{-\ii\varphi_{99}t/2}
  \ket{\tilde\beta\ee^{-\ii\Omega_{99}^{\mathrm{ef}\prime}t/2}}
  \notag\\
  &\qquad\otimes
  \left(
  C\ee^{-\ii\Omega_{99}^{\mathrm{ef}\prime}t/2}\ket{+,9}
  -S\ket{-,9}
  \right)
  \notag\\
  &\quad+
  C\ee^{+\ii\varphi_{99}t/2}
  \ket{\tilde\beta\ee^{+\ii\Omega_{99}^{\mathrm{ef}\prime}t/2}}
  \notag\\
  &\qquad\otimes
  \left(
  S\ee^{+\ii\Omega_{99}^{\mathrm{ef}\prime}t/2}\ket{+,9}
  +C\ket{-,9}
  \right)
  \Bigg].
\end{align}

\subsection{Separaci\'on por ramas polarit\'onicas}

Agrupando los t\'erminos que multiplican a $\ket{+,9}$ y a $\ket{-,9}$,
\begin{equation}
  \boxed{
  \ket{\tilde\Psi(t)}
  =
  \ee^{-\ii A_{99}t}
  \left(
  \ket{\chi_+(t)}\ket{+,9}
  +
  \ket{\chi_-(t)}\ket{-,9}
  \right).
  }
\end{equation}
Las ramas vibracionales condicionadas son
\begin{align}
  \ket{\chi_+(t)}
  &\equiv
  CS
  \Bigg[
  \ee^{\frac{\ii}{2}(\varphi_{99}+\Omega_{99}^{\mathrm{ef}\prime})t}
  \ket{\tilde\beta\ee^{+\ii\Omega_{99}^{\mathrm{ef}\prime}t/2}}
  \notag\\
  &\hspace{2.5cm}
  -
  \ee^{-\frac{\ii}{2}(\varphi_{99}+\Omega_{99}^{\mathrm{ef}\prime})t}
  \ket{\tilde\beta\ee^{-\ii\Omega_{99}^{\mathrm{ef}\prime}t/2}}
  \Bigg],
  \\
  \ket{\chi_-(t)}
  &\equiv
  C^2\ee^{+\ii\varphi_{99}t/2}
  \ket{\tilde\beta\ee^{+\ii\Omega_{99}^{\mathrm{ef}\prime}t/2}}
  \notag\\
  &\quad+
  S^2\ee^{-\ii\varphi_{99}t/2}
  \ket{\tilde\beta\ee^{-\ii\Omega_{99}^{\mathrm{ef}\prime}t/2}}.
\end{align}
Estas ramas no son ortogonales en general. Los estados coherentes forman una
base sobrecompleta, por lo que el traslape entre los dos paquetes vibracionales
es distinto de cero salvo cuando est\'an suficientemente separados en el
espacio de fase.

\section{Traslape entre los paquetes fon\'onicos}

El traslape de dos estados coherentes es
\begin{equation}
  \boxed{
  \braket{\alpha|\beta}
  =
  \exp\left[
  -\frac12\left(|\alpha|^2+|\beta|^2\right)
  +\alpha^*\beta
  \right].
  }
\end{equation}
Para los paquetes que giran en sentidos opuestos se toma
\begin{equation}
  \alpha
  =
  \tilde\beta\ee^{-\ii\Omega_{99}^{\mathrm{ef}\prime}t/2},
  \qquad
  \beta
  =
  \tilde\beta\ee^{+\ii\Omega_{99}^{\mathrm{ef}\prime}t/2}.
\end{equation}
Como $|\alpha|^2=|\beta|^2=M$,
\begin{align}
  \braket{
  \tilde\beta\ee^{-\ii\Omega_{99}^{\mathrm{ef}\prime}t/2}
  |
  \tilde\beta\ee^{+\ii\Omega_{99}^{\mathrm{ef}\prime}t/2}}
  &=
  \exp\left[
  -M+M\ee^{\ii\Omega_{99}^{\mathrm{ef}\prime}t}
  \right]
  \\
  &=
  \exp\left[
  M\left(\cos\Omega_{99}^{\mathrm{ef}\prime}t-1\right)
  +\ii M\sin\Omega_{99}^{\mathrm{ef}\prime}t
  \right].
\end{align}
Usando $\cos x-1=-2\sin^2(x/2)$,
\begin{equation}
  \boxed{
  \braket{
  \tilde\beta\ee^{-\ii\Omega_{99}^{\mathrm{ef}\prime}t/2}
  |
  \tilde\beta\ee^{+\ii\Omega_{99}^{\mathrm{ef}\prime}t/2}}
  =
  \ee^{-2M\sin^2(\Omega_{99}^{\mathrm{ef}\prime}t/2)}
  \ee^{\ii M\sin(\Omega_{99}^{\mathrm{ef}\prime}t)}.
  }
\end{equation}
Para $M=9$,
\begin{equation}
  \boxed{
  \braket{
  \tilde\beta\ee^{-\ii\Omega_{99}^{\mathrm{ef}\prime}t/2}
  |
  \tilde\beta\ee^{+\ii\Omega_{99}^{\mathrm{ef}\prime}t/2}}
  =
  \ee^{-18\sin^2(\Omega_{99}^{\mathrm{ef}\prime}t/2)}
  \ee^{9\ii\sin(\Omega_{99}^{\mathrm{ef}\prime}t)}.
  }
\end{equation}
El factor real exponencial controla la envolvente de la interferencia. Cuando
los paquetes est\'an separados, ese factor se hace peque\~no. La segunda
exponencial es una fase asociada al movimiento de los estados coherentes en el
espacio de fase.

\section{Inversi\'on polarit\'onica en la aproximaci\'on lineal}

La inversi\'on polarit\'onica es
\begin{equation}
  \mathcal W_{\mathrm{lin}}(t)
  =
  \bra{\tilde\Psi(t)}\hat S_z^{(9)}\ket{\tilde\Psi(t)}.
\end{equation}
Como $\hat S_z^{(9)}\ket{\pm,9}=\pm\ket{\pm,9}$,
\begin{equation}
  \boxed{
  \mathcal W_{\mathrm{lin}}(t)
  =
  \braket{\chi_+(t)|\chi_+(t)}
  -
  \braket{\chi_-(t)|\chi_-(t)}.
  }
\end{equation}

Para abreviar, se define
\begin{equation}
  \omega'
  \equiv
  \Omega_{99}^{\mathrm{ef}\prime},
  \qquad
  \Phi
  \equiv
  \varphi_{99}+\omega'.
\end{equation}
Al multiplicar $\bra{\chi_+}$ por $\ket{\chi_+}$ aparecen dos t\'erminos de
norma uno y dos t\'erminos cruzados:
\begin{align}
  \braket{\chi_+|\chi_+}
  &=(CS)^2
  \Bigg[
  2
  -
  \ee^{-\ii\Phi t}
  \braket{
  \tilde\beta\ee^{+\ii\omega't/2}
  |
  \tilde\beta\ee^{-\ii\omega't/2}}
  \notag\\
  &\hspace{2.8cm}
  -
  \ee^{+\ii\Phi t}
  \braket{
  \tilde\beta\ee^{-\ii\omega't/2}
  |
  \tilde\beta\ee^{+\ii\omega't/2}}
  \Bigg].
\end{align}
Los dos t\'erminos cruzados son complejos conjugados. Por tanto,
\begin{equation}
  \boxed{
  \braket{\chi_+|\chi_+}
  =
  2C^2S^2
  \left[
  1-
  \ee^{-2M\sin^2(\omega't/2)}
  \cos\left(
  \Phi t+M\sin\omega't
  \right)
  \right].
  }
\end{equation}

Para la rama inferior,
\begin{align}
  \braket{\chi_-|\chi_-}
  &=
  C^4+S^4
  \notag\\
  &\quad+
  C^2S^2
  \left[
  \ee^{-\ii\varphi_{99}t}
  \braket{
  \tilde\beta\ee^{+\ii\omega't/2}
  |
  \tilde\beta\ee^{-\ii\omega't/2}}
  +\text{c.c.}
  \right].
\end{align}
Entonces
\begin{equation}
  \boxed{
  \braket{\chi_-|\chi_-}
  =
  C^4+S^4
  +
  2C^2S^2
  \ee^{-2M\sin^2(\omega't/2)}
  \cos\left(
  \varphi_{99}t+M\sin\omega't
  \right).
  }
\end{equation}

Sustituyendo ambas normas,
\begin{align}
  \mathcal W_{\mathrm{lin}}(t)
  &=
  2C^2S^2-C^4-S^4
  \notag\\
  &\quad-
  2C^2S^2
  \ee^{-2M\sin^2(\omega't/2)}
  \Big[
  \cos(\varphi_{99}t+M\sin\omega't+\omega't)
  \notag\\
  &\hspace{5.1cm}
  +
  \cos(\varphi_{99}t+M\sin\omega't)
  \Big].
\end{align}
Como
\begin{equation}
  C^4+S^4
  =
  (C^2+S^2)^2-2C^2S^2
  =
  1-2C^2S^2,
\end{equation}
el t\'ermino independiente queda
\begin{equation}
  2C^2S^2-C^4-S^4
  =
  -1+4C^2S^2.
\end{equation}
Se usa adem\'as
\begin{equation}
  \cos(a+b)+\cos a
  =
  2\cos\left(a+\frac b2\right)\cos\frac b2.
\end{equation}
Tomando
\begin{equation}
  a=\varphi_{99}t+M\sin\omega't,
  \qquad
  b=\omega't,
\end{equation}
se obtiene
\begin{align}
  \mathcal W_{\mathrm{lin}}(t)
  &=-1+4C^2S^2
  \Bigg[
  1-
  \ee^{-2M\sin^2(\omega't/2)}
  \notag\\
  &\hspace{2.5cm}\times
  \cos\left(
  \varphi_{99}t+M\sin\omega't+\frac{\omega't}{2}
  \right)
  \cos\frac{\omega't}{2}
  \Bigg].
\end{align}
Finalmente,
\begin{equation}
  4C^2S^2
  =
  1-\frac{\delta_9^2}{(\Omega_{99}^{\mathrm{ef}})^2}
  =
  \frac{(\Omega_{99}^{\mathrm{ef}})^2-\delta_9^2}
  {(\Omega_{99}^{\mathrm{ef}})^2}.
\end{equation}
Por tanto,
\begin{equation}
  \boxed{
  \begin{aligned}
  \mathcal W_{\mathrm{lin}}(t)
  &=-1+
  \frac{(\Omega_{99}^{\mathrm{ef}})^2-\delta_9^2}
  {(\Omega_{99}^{\mathrm{ef}})^2}
  \Bigg[
  1-
  \ee^{-18\sin^2(\Omega_{99}^{\mathrm{ef}\prime}t/2)}
  \\
  &\qquad\times
  \cos\left(
  \varphi_{99}t
  +9\sin(\Omega_{99}^{\mathrm{ef}\prime}t)
  +\frac{\Omega_{99}^{\mathrm{ef}\prime}t}{2}
  \right)
  \cos\left(
  \frac{\Omega_{99}^{\mathrm{ef}\prime}t}{2}
  \right)
  \Bigg].
  \end{aligned}
  }
\end{equation}

\subsection{Caso resonante}

Para el bloque $n=9$, la resonancia ocurre en $\nu=6g$, de modo que
$\delta_9=0$. Entonces
\begin{equation}
  \Omega_{99}^{\mathrm{ef}}=3\nu\lambda,
  \qquad
  \Omega_{99}^{\mathrm{ef}\prime}
  =
  \frac{\nu\lambda}{6},
  \qquad
  \varphi_{99}
  =
  \frac32\nu\lambda.
\end{equation}
La inversi\'on polarit\'onica aproximada se reduce a
\begin{equation}
  \boxed{
  \begin{aligned}
  \mathcal W_{R,\mathrm{lin}}(t)
  &=
  -\exp\left[
  -18\sin^2\left(\frac{\nu\lambda t}{12}\right)
  \right]
  \\
  &\quad\times
  \cos\left[
  \frac32\nu\lambda t
  +9\sin\left(\frac{\nu\lambda t}{6}\right)
  +\frac{\nu\lambda t}{12}
  \right]
  \cos\left(\frac{\nu\lambda t}{12}\right).
  \end{aligned}
  }
\end{equation}

\section{Inversi\'on electr\'onica en la aproximaci\'on lineal}

La inversi\'on electr\'onica se define por
\begin{equation}
  W_{\mathrm{lin}}(t)
  =
  \bra{\tilde\Psi(t)}\hat\sigma_z\ket{\tilde\Psi(t)}.
\end{equation}
En el bloque $n=9$,
\begin{equation}
  \ket{e,8}
  =
  \frac{\ket{+,9}+\ket{-,9}}{\sqrt2},
  \qquad
  \ket{g,9}
  =
  \frac{\ket{+,9}-\ket{-,9}}{\sqrt2}.
\end{equation}
Por ello,
\begin{align}
  \hat\sigma_z
  &=
  \ket{e,8}\bra{e,8}
  -
  \ket{g,9}\bra{g,9}
  \notag\\
  &=
  \boxed{
  \ket{+,9}\bra{-,9}
  +
  \ket{-,9}\bra{+,9}
  }.
\end{align}
La inversi\'on electr\'onica no es una diferencia de normas polarit\'onicas.
En la base polarit\'onica, $\hat\sigma_z$ intercambia las ramas $+$ y $-$.
Por tanto,
\begin{equation}
  W_{\mathrm{lin}}(t)
  =
  \braket{\chi_+|\chi_-}
  +
  \braket{\chi_-|\chi_+}
  =
  \boxed{2\operatorname{Re}\braket{\chi_+|\chi_-}}.
\end{equation}
Al sustituir las ramas y usar el traslape coherente,
\begin{align}
  \operatorname{Re}\braket{\chi_+|\chi_-}
  &=
  CS(C^2-S^2)\cos\left(\frac{\omega't}{2}\right)
  \notag\\
  &\quad+
  CS(S^2-C^2)\ee^{-2M\sin^2(\omega't/2)}
  \cos\left[
  \varphi_{99}t+\omega't+M\sin(\omega't)
  \right].
\end{align}
En consecuencia,
\begin{equation}
  \boxed{
  \begin{aligned}
  W_{\mathrm{lin}}(t)
  &=
  2CS(C^2-S^2)
  \Bigg\{
  \cos\left(
  \frac{\Omega_{99}^{\mathrm{ef}\prime}t}{2}
  \right)
  \\
  &\quad-
  \ee^{-18\sin^2(\Omega_{99}^{\mathrm{ef}\prime}t/2)}
  \cos\left[
  \varphi_{99}t
  +\Omega_{99}^{\mathrm{ef}\prime}t
  +9\sin(\Omega_{99}^{\mathrm{ef}\prime}t)
  \right]
  \Bigg\}.
  \end{aligned}
  }
\end{equation}
En resonancia, $\delta_9=0$ y $C^2=S^2=1/2$. Entonces
\begin{equation}
  \boxed{
  W_{R,\mathrm{lin}}(t)=0.
  }
\end{equation}

\section{Tiempo de revival}

Los paquetes giran con velocidades opuestas
$\pm\Omega_{nM}^{\mathrm{ef}\prime}/2$. Su \'angulo relativo es
\begin{equation}
  \Delta\vartheta(t)=\Omega_{nM}^{\mathrm{ef}\prime}t.
\end{equation}
El primer revival ocurre cuando
$\Omega_{nM}^{\mathrm{ef}\prime}t_r=2\pi$. Por tanto,
\begin{equation}
  \boxed{
  t_r
  =
  \frac{2\pi}{\Omega_{nM}^{\mathrm{ef}\prime}}
  =
  \frac{4\pi\Omega_{nM}^{\mathrm{ef}}}{\nu^2\lambda^2}
  =
  \frac{4\pi}{\nu^2\lambda^2}
  \sqrt{(2g\sqrt n-\nu)^2+\nu^2\lambda^2M}.
  }
\end{equation}
En resonancia,
\begin{equation}
  \boxed{
  t_r^{(R)}=\frac{4\pi\sqrt M}{\nu\lambda}.
  }
\end{equation}
Para $M=9$, $t_r^{(R)}=12\pi/(\nu\lambda)$.

La linealizaci\'on captura los tiempos de revival, pero pierde la curvatura de
$\Omega_{nm}^{\mathrm{ef}}$ dentro del ancho $\Delta m=\sqrt M$. Por eso no
reproduce correctamente el ancho ni la disminuci\'on de altura de los l\'obulos
de revival.

\section{Soluci\'on exacta del hamiltoniano efectivo}

Para recuperar esos efectos se usa la soluci\'on completa en un bloque $n$:
\begin{equation}
  \boxed{
  \begin{aligned}
  \ket{\tilde\Psi(t)}
  &=
  \ee^{-\ii B_nt}\tilde c_0d_n^-\ket{-,n,0}
  \\
  &\quad+
  \sum_{m=1}^{\infty}
  \left[
  p_{nm}^+(t)\ket{+,n,m-1}
  +
  p_{nm}^-(t)\ket{-,n,m}
  \right],
  \end{aligned}
  }
\end{equation}
donde
\begin{equation}
  \boxed{
  B_n=n\Omega+\frac{\nu\lambda^2}{4}-g\sqrt n.
  }
\end{equation}
Para que la notaci\'on de la soluci\'on quede completamente determinada, se
recuerdan las cantidades asociadas con cada doblete polarit\'on--fon\'on:
\begin{equation}
  \boxed{
  \begin{aligned}
  A_{nm}
  &=
  n\Omega+\frac{\nu\lambda^2}{4}
  +\nu\left(m-\frac12\right),
  \\
  \delta_n
  &=2g\sqrt n-\nu,
  \\
  k_m
  &=\nu\lambda\sqrt m,
  \\
  \Omega_{nm}^{\mathrm{ef}}
  &=\sqrt{\delta_n^2+k_m^2}
  =\sqrt{(2g\sqrt n-\nu)^2+\nu^2\lambda^2m}.
  \end{aligned}
  }
\end{equation}
La fase $A_{nm}t$ procede del centro energ\'etico del doblete, mientras que
$\Omega_{nm}^{\mathrm{ef}}$ determina la separaci\'on entre sus dos
eigenenerg\'ias y, por tanto, la frecuencia de intercambio entre las ramas.
Las amplitudes exactas son
\begin{align}
  p_{nm}^+(t)
  &=
  \ee^{-\ii A_{nm}t}
  \Bigg[
  \left(
  \cos\frac{\Omega_{nm}^{\mathrm{ef}}t}{2}
  -
  \ii\frac{\delta_n}{\Omega_{nm}^{\mathrm{ef}}}
  \sin\frac{\Omega_{nm}^{\mathrm{ef}}t}{2}
  \right)\tilde c_{m-1}d_n^+
  \notag\\
  &\hspace{2.8cm}
  +
  \ii\frac{k_m}{\Omega_{nm}^{\mathrm{ef}}}
  \sin\frac{\Omega_{nm}^{\mathrm{ef}}t}{2}
  \tilde c_md_n^-
  \Bigg],
  \\
  p_{nm}^-(t)
  &=
  \ee^{-\ii A_{nm}t}
  \Bigg[
  \ii\frac{k_m}{\Omega_{nm}^{\mathrm{ef}}}
  \sin\frac{\Omega_{nm}^{\mathrm{ef}}t}{2}
  \tilde c_{m-1}d_n^+
  \notag\\
  &\hspace{2.8cm}
  +
  \left(
  \cos\frac{\Omega_{nm}^{\mathrm{ef}}t}{2}
  +
  \ii\frac{\delta_n}{\Omega_{nm}^{\mathrm{ef}}}
  \sin\frac{\Omega_{nm}^{\mathrm{ef}}t}{2}
  \right)\tilde c_md_n^-
  \Bigg].
\end{align}
Para el estado inicial $\ket{-,9}\ket{\tilde\beta}$,
\begin{equation}
  \boxed{
  p_{9m}^+(t)
  =
  \ee^{-\ii A_{9m}t}
  \ii\frac{k_m}{\Omega_{9m}^{\mathrm{ef}}}
  \sin\left(\frac{\Omega_{9m}^{\mathrm{ef}}t}{2}\right)\tilde c_m,
  }
\end{equation}
\begin{equation}
  \boxed{
  p_{9m}^-(t)
  =
  \ee^{-\ii A_{9m}t}
  \left[
  \cos\left(\frac{\Omega_{9m}^{\mathrm{ef}}t}{2}\right)
  +
  \ii\frac{\delta_9}{\Omega_{9m}^{\mathrm{ef}}}
  \sin\left(\frac{\Omega_{9m}^{\mathrm{ef}}t}{2}\right)
  \right]\tilde c_m.
  }
\end{equation}

\section{Inversi\'on polarit\'onica exacta}

Para el bloque $n=9$,
\begin{equation}
  \mathcal W_{\mathrm{ef}}(t)
  =
  -|\tilde c_0|^2
  +
  \sum_{m=1}^{\infty}
  \left(|p_{9m}^+(t)|^2-|p_{9m}^-(t)|^2\right).
\end{equation}
Las probabilidades son
\begin{equation}
  |p_{9m}^+|^2
  =
  |\tilde c_m|^2
  \frac{k_m^2}{(\Omega_{9m}^{\mathrm{ef}})^2}
  \sin^2\left(\frac{\Omega_{9m}^{\mathrm{ef}}t}{2}\right),
\end{equation}
\begin{equation}
  |p_{9m}^-|^2
  =
  |\tilde c_m|^2
  \left[
  \cos^2\left(\frac{\Omega_{9m}^{\mathrm{ef}}t}{2}\right)
  +
  \frac{\delta_9^2}{(\Omega_{9m}^{\mathrm{ef}})^2}
  \sin^2\left(\frac{\Omega_{9m}^{\mathrm{ef}}t}{2}\right)
  \right].
\end{equation}
Su diferencia puede escribirse como
\begin{align}
  |p_{9m}^+|^2-|p_{9m}^-|^2
  &=
  -|\tilde c_m|^2
  \notag\\
  &\quad+
  |\tilde c_m|^2
  \frac{k_m^2}{(\Omega_{9m}^{\mathrm{ef}})^2}
  [1-\cos(\Omega_{9m}^{\mathrm{ef}}t)].
\end{align}
Usando la normalizaci\'on de la distribuci\'on fon\'onica,
\begin{equation}
  \boxed{
  \mathcal W_{\mathrm{ef}}(t)
  =
  -1+
  \sum_{m=1}^{\infty}
  |\tilde c_m|^2
  \frac{k_m^2}{(\Omega_{9m}^{\mathrm{ef}})^2}
  [1-\cos(\Omega_{9m}^{\mathrm{ef}}t)].
  }
\end{equation}

\subsection{Caso resonante}

En resonancia, $\Omega_{9m}^{\mathrm{ef}}=k_m=\nu\lambda\sqrt m$, y
\begin{equation}
  \boxed{
  \mathcal W_{R,\mathrm{ef}}(t)
  =
  -
  \sum_{m=0}^{\infty}
  |\tilde c_m|^2\cos(\nu\lambda\sqrt m\,t).
  }
\end{equation}
Para $\tilde\beta=3$,
\begin{equation}
  \boxed{
  \mathcal W_{R,\mathrm{ef}}(t)
  =
  -\ee^{-9}
  \sum_{m=0}^{\infty}
  \frac{9^m}{m!}\cos(\nu\lambda\sqrt m\,t).
  }
\end{equation}
El desfasaje de las frecuencias proporcionales a $\sqrt m$ produce el colapso;
su refase parcial produce revivals con anchura y altura variables.

\section{Inversi\'on electr\'onica exacta}

La inversi\'on electr\'onica requiere un cuidado adicional porque el operador
$\hat\sigma_z$ es no diagonal en la base polarit\'onica. Para cualquier n\'umero
fon\'onico $m$ se tiene
\begin{equation}
  \hat\sigma_z
  =
  \ket{+,9,m}\bra{-,9,m}
  +
  \ket{-,9,m}\bra{+,9,m}.
\end{equation}
Por tanto, solamente interfieren amplitudes que tengan el mismo n\'umero de
fonones. En la expansi\'on exacta, $p_{9m}^-$ multiplica a
$\ket{-,9,m}$, pero $p_{9m}^+$ multiplica a $\ket{+,9,m-1}$. Para obtener la
amplitud de la rama superior con $m$ fonones debe reemplazarse $m$ por $m+1$.
De esta observaci\'on se sigue
\begin{equation}
  \boxed{
  W_{\mathrm{ef}}(t)
  =
  2\operatorname{Re}
  \sum_{m=0}^{\infty}
  p_{9,m+1}^{+*}(t)p_{9m}^-(t),
  }
\end{equation}
donde, para incluir el estado aislado de $m=0$, se entiende
\begin{equation}
  p_{90}^-(t)=\ee^{-\ii B_9t}\tilde c_0.
\end{equation}

Para hacer transparente el producto de amplitudes, se introducen las
abreviaciones
\begin{equation}
  \mathcal C_m(t)
  =
  \cos\left(\frac{\Omega_{9m}^{\mathrm{ef}}t}{2}\right),
  \qquad
  \mathcal S_m(t)
  =
  \sin\left(\frac{\Omega_{9m}^{\mathrm{ef}}t}{2}\right),
\end{equation}
\begin{equation}
  D_m=\frac{\delta_9}{\Omega_{9m}^{\mathrm{ef}}},
  \qquad
  K_m=\frac{k_m}{\Omega_{9m}^{\mathrm{ef}}}.
\end{equation}
Para $m\geq1$, las amplitudes toman la forma
\begin{equation}
  p_{9,m+1}^{+*}(t)
  =
  -\ii\ee^{+\ii A_{9,m+1}t}
  K_{m+1}\mathcal S_{m+1}(t)\tilde c_{m+1}^*,
\end{equation}
\begin{equation}
  p_{9m}^-(t)
  =
  \ee^{-\ii A_{9m}t}
  [\mathcal C_m(t)+\ii D_m\mathcal S_m(t)]\tilde c_m.
\end{equation}
Al multiplicarlas,
\begin{align}
  p_{9,m+1}^{+*}p_{9m}^-
  &=
  \ee^{\ii(A_{9,m+1}-A_{9m})t}
  K_{m+1}\mathcal S_{m+1}
  \tilde c_{m+1}^*\tilde c_m
  \notag\\
  &\quad\times
  [D_m\mathcal S_m-\ii\mathcal C_m].
\end{align}
La diferencia entre los centros energ\'eticos de dobletes consecutivos es
\begin{align}
  A_{9,m+1}-A_{9m}
  &=
  \nu\left(m+1-\frac12\right)
  -\nu\left(m-\frac12\right)
  \notag\\
  &=\nu.
\end{align}
Si $\tilde\beta$ es real, como ocurre para $\tilde\beta=3$, los coeficientes
$\tilde c_m$ tambi\'en son reales. Usando
\begin{equation}
  \operatorname{Re}
  \left\{
  \ee^{\ii\nu t}
  [D_m\mathcal S_m-\ii\mathcal C_m]
  \right\}
  =
  D_m\mathcal S_m\cos(\nu t)
  +\mathcal C_m\sin(\nu t),
\end{equation}
se obtiene, para todos los t\'erminos con $m\geq1$,
\begin{align}
  2\operatorname{Re}
  [p_{9,m+1}^{+*}p_{9m}^-]
  &={}
  2K_{m+1}\mathcal S_{m+1}	ilde c_{m+1}\tilde c_m
  \notag\\
  &\quad\times
  [D_m\mathcal S_m\cos(\nu t)
  +\mathcal C_m\sin(\nu t)].
\end{align}

El t\'ermino de borde $m=0$ debe calcularse por separado, puesto que
$\ket{-,9,0}$ no pertenece a ning\'un doblete. En este caso,
\begin{align}
  p_{91}^{+*}p_{90}^-
  &={}
  -\ii
  \ee^{\ii(A_{91}-B_9)t}
  K_1\mathcal S_1\tilde c_1\tilde c_0,
  \\
  2\operatorname{Re}[p_{91}^{+*}p_{90}^-]
  &={}
  2\sin[(A_{91}-B_9)t]
  K_1\mathcal S_1\tilde c_1\tilde c_0.
\end{align}
Sumando el t\'ermino de borde y todos los dobletes se llega a
\begin{equation}
  \boxed{
  \begin{aligned}
  W_{\mathrm{ef}}(t)
  &={}
  2\sin[(A_{91}-B_9)t]
  K_1\mathcal S_1(t)\tilde c_1\tilde c_0
  \\
  &\quad+
  2\sum_{m=1}^{\infty}
  K_{m+1}\mathcal S_{m+1}(t)
  \tilde c_{m+1}\tilde c_m
  \\
  &\qquad\times
  [D_m\mathcal S_m(t)\cos(\nu t)
  +\mathcal C_m(t)\sin(\nu t)].
  \end{aligned}
  }
\end{equation}
Esta expresi\'on muestra dos escalas din\'amicas. Las funciones
$\mathcal C_m$ y $\mathcal S_m$ contienen las frecuencias de Rabi efectivas
dependientes de $m$, mientras que $\sin(\nu t)$ y $\cos(\nu t)$ corresponden
a la fase relativa entre dobletes fon\'onicos consecutivos.

\subsection{Caso resonante}

En resonancia polarit\'on--fon\'on, $\delta_9=0$. Por ello,
\begin{equation}
  D_m=0,
  \qquad
  K_m=1,
  \qquad
  \Omega_{9m}^{\mathrm{ef}}=\nu\lambda\sqrt m.
\end{equation}
Adem\'as, $A_{91}-B_9=\nu$ cuando $\nu=6g$. El t\'ermino de borde puede
incorporarse a la suma si se usan
$\mathcal C_0=1$ y
$\mathcal S_1=\sin(\nu\lambda t/2)$. El resultado se reduce a
\begin{equation}
  \boxed{
  \begin{aligned}
  W_{R,\mathrm{ef}}(t)
  ={}&
  2\sum_{m=0}^{\infty}
  \sin(\nu t)
  \sin\left(\frac{\nu\lambda\sqrt{m+1}\,t}{2}\right)
  \\
  &\qquad\times
  \cos\left(\frac{\nu\lambda\sqrt m\,t}{2}\right)
  \tilde c_{m+1}\tilde c_m.
  \end{aligned}
  }
\end{equation}
Aunque la inversi\'on electr\'onica de la aproximaci\'on lineal se anula en
resonancia, la soluci\'on exacta contiene interferencias entre sectores
fon\'onicos vecinos. Estas interferencias conservan una din\'amica r\'apida a
frecuencia $\nu$, modulada por el colapso y los revivals asociados con la
distribuci\'on coherente.

\section{Resultados gr\'aficos}

Las siguientes figuras comparan la aproximaci\'on lineal, la soluci\'on exacta
del hamiltoniano efectivo y la evoluci\'on num\'erica del hamiltoniano completo
de Holstein--Cummings. En todas las comparaciones se usa
$\tilde\beta=3$, de modo que $M=9$.

\subsection{Aproximaci\'on lineal de la frecuencia efectiva}

\begin{figure}[H]
  \centering
  \includegraphics[width=0.95\linewidth]{\detokenize{../../notas-doctorado/imgs/Estado coherente fononico en bloque polaritonico/InversionP_aprox_lineal.png}}
  \caption{Inversi\'on polarit\'onica obtenida con la aproximaci\'on lineal de
  la frecuencia de Rabi efectiva. La curva amarilla corresponde al caso
  resonante, $\nu=6g$, y la azul a $\nu=5.99g$. La linealizaci\'on predice
  correctamente los tiempos de revival, determinados por la reuni\'on de los
  dos paquetes coherentes, pero no reproduce la anchura ni la reducci\'on de
  altura de los l\'obulos de revival porque elimina la curvatura de
  $\Omega_{9m}^{\mathrm{ef}}$ respecto de $m$.}
\end{figure}

\begin{figure}[H]
  \centering
  \includegraphics[width=0.95\linewidth]{\detokenize{../../notas-doctorado/imgs/Estado coherente fononico en bloque polaritonico/InversionE_aprox_lineal.png}}
  \caption{Inversi\'on electr\'onica dentro de la aproximaci\'on lineal. En
  resonancia, curva amarilla, el resultado es id\'enticamente cero porque
  $C^2=S^2=1/2$: cada polarit\'on contiene siempre probabilidades iguales para
  los estados electr\'onicos excitado y base. Fuera de resonancia, curva azul
  para $\nu=5.99g$, la mezcla deja de ser exactamente sim\'etrica y aparece una
  se\~nal distinta de cero; sin embargo, la linealizaci\'on todav\'ia omite la
  din\'amica producida por la curvatura de las frecuencias efectivas.}
\end{figure}

\subsection{Soluci\'on exacta del hamiltoniano efectivo}

\begin{figure}[H]
  \centering
  \includegraphics[width=0.95\linewidth]{\detokenize{../../notas-doctorado/imgs/Estado coherente fononico en bloque polaritonico/InversionP.png}}
  \caption{Inversi\'on polarit\'onica calculada con la dependencia exacta de
  $\Omega_{9m}^{\mathrm{ef}}$. La curva amarilla corresponde a resonancia y la
  azul al caso fuera de resonancia. Ahora aparecen los colapsos y revivals
  caracter\'isticos del modelo JCP, junto con el ensanchamiento y la reducci\'on
  de altura de sus l\'obulos. La desintonizaci\'on introduce adem\'as el t\'ermino
  $\delta_9^2/(\Omega_{9m}^{\mathrm{ef}})^2$, que impide una transferencia
  polarit\'onica completa y desplaza verticalmente la se\~nal.}
\end{figure}

\begin{figure}[H]
  \centering
  \includegraphics[width=0.95\linewidth]{\detokenize{../../notas-doctorado/imgs/Estado coherente fononico en bloque polaritonico/InversionE.png}}
  \caption{Inversi\'on electr\'onica del hamiltoniano efectivo: caso resonante
  en amarillo y no resonante en azul. Fuera de resonancia, el cociente
  $D_m=\delta_9/\Omega_{9m}^{\mathrm{ef}}$ deja de ser cero y la observable
  combina las fases de sectores vecinos $m$ y $m+1$ con la oscilaci\'on r\'apida
  de frecuencia $\nu$. El batimiento resultante desplaza los m\'aximos de
  refase, por lo que los grupos de oscilaciones azules parecen acercarse entre
  s\'i a lo largo del intervalo mostrado.}
\end{figure}

\subsection{Comparaci\'on con el hamiltoniano completo}

\begin{figure}[H]
  \centering
  \includegraphics[width=0.95\linewidth]{\detokenize{../../notas-doctorado/imgs/Estado coherente fononico en bloque polaritonico/ComparacionP.png}}
  \caption{Comparaci\'on, en resonancia, de la inversi\'on polarit\'onica
  generada por el hamiltoniano efectivo y por el hamiltoniano completo de
  Holstein--Cummings resuelto num\'ericamente. Las curvas son pr\'acticamente
  indistinguibles, lo que muestra que el modelo efectivo reproduce de manera
  precisa la din\'amica polarit\'onica completa en el r\'egimen y la escala
  temporal considerados.}
\end{figure}

\begin{figure}[H]
  \centering
  \includegraphics[width=0.95\linewidth]{\detokenize{../../notas-doctorado/imgs/Estado coherente fononico en bloque polaritonico/ComparacionE.png}}
  \caption{Comparaci\'on de las inversiones electr\'onicas en resonancia. La
  evoluci\'on obtenida con el hamiltoniano efectivo coincide con la del modelo
  completo, incluso para la modulaci\'on r\'apida debida a la frecuencia
  vibracional y para la envolvente de colapsos y revivals.}
\end{figure}

\begin{figure}[H]
  \centering
  \includegraphics[width=0.95\linewidth]{\detokenize{../../notas-doctorado/imgs/Estado coherente fononico en bloque polaritonico/Fidelidad.png}}
  \caption{Fidelidad entre los estados evolucionados con el hamiltoniano
  efectivo y con el hamiltoniano completo de Holstein--Cummings. Su valor se
  mantiene pr\'acticamente igual a uno durante el intervalo estudiado. Esta
  proximidad a la unidad explica la coincidencia de las inversiones
  polarit\'onica y electr\'onica en las dos figuras anteriores.}
\end{figure}

\section{Lectura f\'isica final}

El estado coherente fon\'onico convierte un solo bloque polarit\'onico en una
superposici\'on de dobletes con diferentes n\'umeros vibracionales. La
dependencia $\Omega_{9m}^{\mathrm{ef}}\propto\sqrt m$ hace que esas componentes
se desfasen inicialmente, produciendo el colapso, y vuelvan a alinearse de
manera parcial en los tiempos de revival. La aproximaci\'on lineal conserva la
separaci\'on angular media entre los paquetes y, por ello, determina bien esos
tiempos; la soluci\'on exacta conserva tambi\'en la curvatura de la dispersi\'on
y recupera la forma completa de cada l\'obulo.

La comparaci\'on con el hamiltoniano completo muestra, finalmente, que la
reducci\'on al hamiltoniano efectivo no solamente conserva una escala de
frecuencia, sino tambi\'en el estado cu\'antico y las interferencias necesarias
para describir simult\'aneamente las inversiones polarit\'onica y electr\'onica
durante el intervalo analizado.

\end{document}
